% ===================================================================
%  BẢNG Số 7 — Làng quê mùa đông. Ngôi nhà quê mùa xuân.
%
%  Traduction, faite en 2026, en vietnamien standard d'aujourd'hui,
%  sur l'ido de texto/io. Regles de la colonne : voir 10-bang-01.tex.
%  Deux scenes, deux numerotations : les cles portent c1 et c2.
%
%  LA PREMIERE SCENE EST UNE SCENE DE NEIGE, ET LE VIETNAMIEN N'EN A
%  PAS. C'est la difficulte du tableau, et elle ne se resout pas par
%  un emprunt : « tuyết » existe, il est chinois, il est ancien dans
%  la langue ecrite, et tout lecteur le comprend — mais il ne traine
%  derriere lui aucun des mots que le francais a autour de la neige.
%  Le bonhomme de neige devient « người tuyết », l'homme de neige ; la
%  bataille de boules de neige, « ném bóng tuyết », jeter des boules
%  de neige ; la glissoire, « chỗ trượt », l'endroit ou l'on glisse ;
%  le traineau, « xe trượt tuyết », la voiture qui glisse sur la
%  neige. Aucun de ces mots n'est un tour de force : ce sont des
%  composes transparents, et c'est ainsi qu'une langue accueille un
%  climat qu'elle n'a pas.
%
%  LA SECONDE SCENE, ELLE, EST CHEZ ELLE. La cour de ferme de mai est
%  le passage du livret ou le vietnamien a le plus a offrir, et le
%  plus a distinguer : la ou le francais dit « bœuf », « vache »,
%  « veau », « taureau », il dit tous les quatre avec le meme mot de
%  base, « bò », precise d'un second : bò cái la vache, bò đực le
%  taureau, bê le veau. Et il separe ce que le francais confond —
%  « lợn » le porc, « lợn nái » la truie, « lợn con » le porcelet ;
%  « gà » la poule, « gà trống » le coq, « gà con » le poussin.
%
%  LE RIZ N'EST PAS LA, ET C'EST VOULU. La ferme de la planche fauche
%  du foin et laboure un champ de ble ; « lúa mì », le ble, s'ecrit
%  litteralement « riz de froment », parce que le riz est ici la
%  cereale par defaut. On garde le compose : c'est ainsi que la langue
%  a nomme la cereale des autres.
% ===================================================================

%%K t07-noto-1 noto
\VUnotes{143.2mm}{%
(*) Về chi tiết các bữa ăn, xem bảng 6 và 13.%
}

%%K t07-apar-1 apar
\VUsaut{4.0mm}
\VUcentre{13.2pt}{90}{LOẠT THỨ HAI}
\VUsaut{3.0mm}
\VUfilet{16mm}
\VUsaut{5.6mm}
\VUtitre{15.0pt}{60}{BẢNG Số 7}
\VUsaut{3.0mm}

%%K t07-tit-1 sub
\VUcentre{12.0pt}{0}{\textit{Cảnh thứ nhất.}}
\VUsaut{3.0mm}

%%K t07-tit-2 sub
\VUpk{10.2pt}{40}{Làng quê mùa đông.}
\VUsaut{4.0mm}

%%K t07-c1-01-1 p
1. --- Trong kỳ nghỉ đầu năm, em theo cha đến Tân Thị, một làng nhỏ nơi cha có mấy việc
buôn bán phải giải quyết.

%%K t07-c1-02-1 p
2. --- Chúng em đi \VUgras{xe khách} \textsuperscript{(11)} chạy giữa thành phố và thị
trấn ấy; nhưng vì trời rất lạnh nên chuyến đi trong chiếc xe ít tiện nghi không mấy dễ
chịu.

%%K t07-c1-03-1 p
3. --- Vì thế chúng em thật sự vui khi thấy \VUgras{người đánh xe} dừng xe trên quảng
trường, trước \VUgras{quán trọ} \textsuperscript{(2)} \textit{Gấu Trắng}. \VUgras{Chủ
quán} \textsuperscript{(4)} đang đợi chúng em trên thềm nhà, bình thản hút \VUgras{tẩu}
của mình.

%%K t07-c1-03-2 p
Ông mời chúng em vào và em chẳng để phải nài, đến ngồi ngay trước ngọn lửa cành nho nổ
lách tách trong lò. Lúc ấy em thật sự chế nhạo \VUgras{ngọn gió bấc} buốt giá làm kẽo
kẹt tấm \VUgras{biển hiệu} \textsuperscript{(3)} cũ và cuốn \VUgras{làn khói}
\textsuperscript{(1)} từ ống khói ra thành những xoáy đen.

%%K t07-c1-04-1 p
4. --- Đã đến giờ ăn trưa. Không chờ lâu hơn, chúng em ăn ngon lành bữa cơm đơn sơ mà
đậm đà người ta dọn cho (*).

%%K t07-tit-3 sub
\VUpk{10.2pt}{40}{Mấy cuộc thăm viếng.}
\VUsaut{3.0mm}

%%K t07-c1-05-1 p
5. --- Sau bữa trưa, em theo cha, người làm nghề bảo hiểm, đi thăm khách hàng. Trước hết
chúng em đến nhà \VUgras{bác thợ rèn} \textsuperscript{(5)}, ở gần quán trọ. Bác đang
bận đóng móng cho một con ngựa. Em thán phục sức khỏe của người phụ việc, người giữ chắc
trên chiếc tạp dề da của mình cái \VUgras{móng} ngựa, trong khi bác thợ rèn đóng chặt
\VUgras{miếng sắt} \textsuperscript{(6)} bằng những nhát búa mạnh.

%%K t07-c1-06-1 p
6. --- Rồi chúng em đến nhà \VUgras{bác thợ giày} \textsuperscript{(7)} của làng, cụ
Giacôbê. Tuy có tấm biển là một chiếc \VUgras{ủng} rất đẹp, cụ chỉ đang đóng lại đế cho
một đôi giày cũ thủng.

%%K t07-c1-06-2 p
Cụ cười bảo chúng em : « Ở thành phố, tôi là "thợ đóng giày" và tôi chẳng có khách. Ở
đây, tôi là "thợ chữa giày" và việc thì không hết. »

%%K t07-c1-07-1 p
7. --- Cụ nói với chúng em về ông hàng xóm là \VUgras{bác thợ cạo}
\textsuperscript{(8)}, người mà khách không hài lòng. \VUgras{Dao cạo} của bác lúc nào
cũng mẻ và \VUgras{kéo} của bác chẳng bao giờ cắt cho ra hồn. Nhưng vì bác là người thợ
cạo duy nhất trong làng nên ai cũng đành để bác cạo mặt và cắt tóc cho. Vả lại đó là một
người keo kiệt và khó chịu.

%%K t07-c1-08-1 p
8. --- May thay, những \VUgras{người hàng xóm} khác không giống bác. Chẳng hạn
\VUgras{bác thợ đóng xe} \textsuperscript{(9)} là một người tốt bụng, chăm chỉ và sẵn
lòng giúp đỡ. Đưa xe cho bác chữa là một điều thú vị. Việc bác làm bao giờ cũng xong
xuôi chu đáo.

%%K t07-c1-09-1 p
9. --- Cụ Giacôbê cũng không phàn nàn gì về \VUgras{bác thợ làm yên}
\textsuperscript{(10)}, người mà đôi khi cụ giúp khâu \VUgras{bộ yên cương} khi công
việc gấp.

%%K t07-c1-09-2 p
Cha em hỏi thăm gia đình bác : « Mọi người đều khỏe. Cháu gái tôi đi \VUgras{học}
\textsuperscript{(12)}. \VUgras{Cô giáo} \textsuperscript{(13)} rất hài lòng về cháu,
cháu là một \VUgras{học sinh} \textsuperscript{(14)} giỏi. Cháu không giống thằng con
nghịch ngợm của tôi; nó chỉ nghĩ đến chơi. \VUgras{Thầy giáo}
\textsuperscript{(15)} chẳng dạy nổi nó điều gì. »

%%K t07-tit-4 sub
\VUsaut{3.0mm}
\VUpk{10.2pt}{40}{Chuyện làng.}
\VUsaut{3.0mm}

%%K t07-c1-10-1 p
10. --- Cụ kể tiếp cho chúng em nghe tin tức trong làng. Người ta sắp xây một trường học
mới, vì trường đặt trong \VUgras{trụ sở làng} đã trở nên quá chật. Vả lại việc ấy dễ
thôi, vì chỉ cần mua một phần mảnh vườn của \VUgras{ông thợ xay}. Đó sẽ là món lợi cho
con người đáng thương ấy. Vì trời lạnh, \VUgras{cối xay} của \VUgras{nhà máy xay}
\textsuperscript{(16)} không xay được lúa mì nữa, bởi \VUgras{bánh xe}
\textsuperscript{(17)} đã bị \VUgras{băng} \textsuperscript{(25)} của \VUgras{con suối}
\textsuperscript{(23)} làm cho đứng lại.

%%K t07-c1-11-1 p
11. --- « Nói đến xây cất, các ông đã xem \VUgras{nhà thờ} \textsuperscript{(18)} mới
của chúng tôi chưa? Nó đẹp lắm dưới lớp áo tuyết. Những con \VUgras{quạ}
\textsuperscript{(19)}, không còn nhận ra tháp chuông cũ của mình, buồn bã đậu trên cây
lớn của \VUgras{nghĩa địa} \textsuperscript{(20)}. »

%%K t07-c1-12-1 p
12. --- Ý nghĩ về nghĩa địa nhắc bác thợ giày nhớ đến những người cha em từng quen và đã
mất năm ngoái. « Bao nhiêu là \VUgras{ngôi mộ} \textsuperscript{(21)}, thưa ông, đã thêm
vào những ngôi mộ khác, dưới hàng \VUgras{trắc bá} \textsuperscript{(22)}! Cái chết đã
cướp đi ông chưởng khế già của chúng tôi. Cả làng đã dự \VUgras{đám tang} của ông. »

%%K t07-c1-13-1 p
13. --- « Kìa người kế nhiệm ông đang trượt băng trên con suối. Ông ấy vừa mang đến nhờ
tôi chữa cái quai \VUgras{giày trượt} \textsuperscript{(26)} bị đứt. »

%%K t07-c1-14-1 p
14. --- « Kìa nữa, trên bờ suối, là \VUgras{người canh đồng} \textsuperscript{(27)} mới.
Đó là một cựu hiến binh làm lũ trẻ nghịch trong làng khiếp sợ. Chúng chạy trốn ngay khi
trông thấy đôi \VUgras{ghệt} \textsuperscript{(29)} và chiếc \VUgras{mũ lưỡi trai cứng}
\textsuperscript{(28)} của ông. »

%%K t07-c1-15-1 p
15. --- « À! kìa cụ Giuse dắt một \VUgras{xe củi bó} \textsuperscript{(30)} đến cho ông
hàng bánh. --- Đúng thế, cha em nói, tôi nhận ra cỗ \VUgras{la} \textsuperscript{(31)}
kéo xe của cụ. Tôi đang cần nói chuyện với cụ, tôi sẽ nhân dịp này. Chào cụ Giacôbê. »

%%K t07-c1-16-1 p
16. --- Vừa lúc chúng em bước ra, trẻ con trong làng tan học và chạy ào tới \VUgras{chỗ
trượt} \textsuperscript{(33)}, liều ngã và gãy tay chân trong cú \VUgras{ngã}
\textsuperscript{(34)}. Một đứa lấy chiếc \VUgras{xe trượt tuyết}
\textsuperscript{(32)} của mình ở nhà bác thợ đóng xe, đặt một bạn lên đó rồi chạy đi.

%%K t07-c1-17-1 p
17. --- Những đứa khác lao tới một \VUgras{đống tuyết} \textsuperscript{(37)} gần
\VUgras{cây cầu} \textsuperscript{(40)} và bắt đầu ném vào \VUgras{người tuyết}
\textsuperscript{(39)} mà chúng đã đắp hôm qua và đội cho cái mũ, cắm cái tẩu, đeo chiếc
cặp sách chéo vai.

%%K t07-c1-18-1 p
18. --- Chúng chẳng mấy bận tâm đến gió buốt và cái lạnh. Phần lớn thậm chí không có
\VUgras{khăn quàng cổ} \textsuperscript{(35)}, và để sưởi ấm sau trận \VUgras{ném bóng
tuyết} \textsuperscript{(38)}, chúng chỉ cần một nắm \VUgras{hạt dẻ}
\textsuperscript{(42)} thật nóng mua của bà \VUgras{hàng} Giacôbina già, người đang run
lên vì lạnh dưới tấm \VUgras{khăn choàng} \textsuperscript{(43)} quá mỏng.

%%K t07-tit-5 sub
\VUsaut{3.0mm}
\VUcentre{12.0pt}{0}{\textit{Cảnh thứ hai.}}
\VUsaut{3.0mm}

%%K t07-tit-6 sub
\VUpk{10.2pt}{40}{Ngôi nhà quê mùa xuân.}
\VUsaut{4.0mm}

%%K t07-c2-01-1 p
1. --- Dù mùa đông có bao thú vui, chúng ta vẫn vui mừng sâu xa đón \VUgras{mùa xuân}
trở lại.

%%K t07-c2-01-2 p
Một cái sân trại, vào buổi chiều tháng năm, mới là một bức tranh vui mắt làm sao!

%%K t07-c2-02-1 p
2. --- Ở chân trời, \VUgras{mặt trời} \textsuperscript{(1)} từ từ lặn, tràn ngập
\VUgras{đồng quê} \textsuperscript{(3)} bằng những \VUgras{tia nắng}
\textsuperscript{(2)} đỏ tía. \VUgras{Người thợ cày} \textsuperscript{(6)} cần mẫn vội
cày nốt \VUgras{thửa ruộng} \textsuperscript{(4)} của mình.

%%K t07-c2-02-2 p
Với chiếc \VUgras{cày} \textsuperscript{(5)} do một con \VUgras{ngựa}
\textsuperscript{(7)} khỏe kéo, ông vạch \VUgras{luống cày} \textsuperscript{(8)} mà vào
đó \VUgras{người gieo hạt} \textsuperscript{(9)} sẽ vãi từng nắm \VUgras{hạt giống} để
sinh ra những bông lúa vàng.

%%K t07-c2-03-1 p
3. --- \VUgras{Thợ cắt cỏ} \textsuperscript{(11)}, với những chiếc hái của mình, đã cắt
cỏ đồng; thợ phơi cỏ đã phơi khô \VUgras{cỏ khô} \textsuperscript{(10)} bằng cách trở nó
với \VUgras{chĩa} \textsuperscript{(12)} và cào, và kìa họ đang trở về.

%%K t07-c2-03-2 p
Người thì đi trước chiếc xe nặng do bò kéo; họ vác dụng cụ trên vai. Người khác, lười
hơn, đã ngồi thoải mái trên đống cỏ thơm.

%%K t07-c2-04-1 p
4. --- Đi qua, họ chào thân mật \VUgras{người làm vườn} \textsuperscript{(16)} đang bận
trong \VUgras{vườn quả} \textsuperscript{(13)}, tỉa \VUgras{cây ăn quả} và ghép
\VUgras{cây lê} \textsuperscript{(14)}. \VUgras{Cây mận} \textsuperscript{(15)} đang nở
hoa, và trong \VUgras{vườn rau} \textsuperscript{(17)} bón kỹ, rau, cải bắp và xà lách
đã tươi tốt.

%%K t07-c2-04-2 p
Trên bức tường thấp, một con \VUgras{công} \textsuperscript{(18)} đẹp xòe đuôi để khoe
những chiếc lông tuyệt đẹp của mình.

%%K t07-tit-7 sub
\VUpk{10.2pt}{40}{Sân trại.}
\VUsaut{3.0mm}

%%K t07-c2-05-1 p
5. --- Cửa \VUgras{chuồng ngựa} \textsuperscript{(19)} mở, và ta thấy máng cỏ cùng
\VUgras{ổ rơm} \textsuperscript{(23)} của con \VUgras{ngựa cái} \textsuperscript{(21)}
mà \VUgras{người coi ngựa} \textsuperscript{(20)} vừa tháo khỏi xe. Con \VUgras{ngựa
con} \textsuperscript{(22)}, trông thật ngộ với đôi chân dài, tung tăng theo mẹ.

%%K t07-c2-06-1 p
6. --- Gần \VUgras{cái giếng} \textsuperscript{(29)}, những con \VUgras{bò cái}
\textsuperscript{(25)} đến uống trong chiếc \VUgras{máng} \textsuperscript{(26)} gỗ dùng
làm chỗ cho chúng uống. Con \VUgras{bê} \textsuperscript{(24)} nhân dịp bú mẹ, còn con
\VUgras{bò đực} \textsuperscript{(27)}, ngửi thấy mùi cỏ khô thơm, rống lên vui vẻ và
kiêu hãnh ngẩng cái đầu với đôi \VUgras{sừng} \textsuperscript{(28)} lợi hại.

%%K t07-c2-07-1 p
7. --- Ai nấy đều làm việc. \VUgras{Chị giúp việc} \textsuperscript{(32)} cầm trong tay
\VUgras{sợi dây} \textsuperscript{(31)} của giếng. Chị múc nước bằng \VUgras{cái gàu}
\textsuperscript{(30)} để cho gia súc uống. Bên cạnh \VUgras{nhà kho}
\textsuperscript{(33)}, một người đàn ông cào rơm, và trước cửa \VUgras{ngôi nhà}
\textsuperscript{(34)}, một bà cụ \VUgras{kéo sợi} \textsuperscript{(35)}, với guồng
quay và \VUgras{con cúi}, đang xe \VUgras{sợi lanh} hay \VUgras{sợi gai} như thuở xưa.

%%K t07-tit-8 sub
\VUsaut{3.0mm}
\VUpk{10.2pt}{40}{Người chủ trại.}
\VUsaut{3.0mm}

%%K t07-c2-08-1 p
8. --- \VUgras{Người chủ trại} \textsuperscript{(36)} chỉ huy tất cả những việc ấy và
dặn dò người làm. Ông nhắc phải cất \VUgras{cái bừa} \textsuperscript{(40)} và cái
\VUgras{cuốc chim} \textsuperscript{(39)} bỏ quên gần \VUgras{cũi}
\textsuperscript{(38)} của con \VUgras{chó} \textsuperscript{(37)}.

%%K t07-c2-09-1 p
9. --- Tiếng ông làm một con \VUgras{bồ câu} \textsuperscript{(41)} hoảng sợ bay vụt về
\VUgras{chuồng bồ câu} \textsuperscript{(42)}, và những con \VUgras{gà mái}
\textsuperscript{(43)} vội chạy về \VUgras{chuồng gà} \textsuperscript{(44)}.

%%K t07-c2-10-1 p
10. --- Trước \VUgras{chuồng cừu} \textsuperscript{(48)}, kìa cụ \VUgras{chăn cừu}
\textsuperscript{(45)} già đang lùa \VUgras{đàn} \textsuperscript{(49)} về. Tay cầm cây
\VUgras{gậy chăn cừu} \textsuperscript{(46)}, thứ mà cụ không bao giờ rời cũng như chiếc
\VUgras{áo choàng} \textsuperscript{(47)} bạc màu của mình, cụ dẫn về chuồng những con
\VUgras{cừu đực} \textsuperscript{(50)}, \VUgras{cừu cái} \textsuperscript{(53)},
\VUgras{cừu con} \textsuperscript{(54)}, \VUgras{dê cái} \textsuperscript{(51)} và
\VUgras{dê con} \textsuperscript{(52)}. Con \VUgras{chó} \textsuperscript{(55)} trung
thành của cụ, tên Argus, đi sau cùng và giục những con chậm chân bằng tiếng sủa.

%%K t07-c2-11-1 p
11. --- Tất cả tiếng động và sự nhộn nhịp ấy không làm mất vẻ thản nhiên của con
\VUgras{bò sữa} \textsuperscript{(56)} hiền lành đang rung chiếc \VUgras{lục lạc}
\textsuperscript{(57)} trong lúc người ta vắt sữa nó trước cửa \VUgras{chuồng bò}
\textsuperscript{(58)}.

%%K t07-c2-12-1 p
12. Nó có vẻ hiểu rằng nó phải cho sữa để \VUgras{bà chủ trại} \textsuperscript{(60)}
làm được \VUgras{bơ} trong chiếc \VUgras{thùng đánh bơ} \textsuperscript{(61)} hay những
\VUgras{tảng phô mai} \textsuperscript{(64)} tuyệt hảo đang nhỏ nước trên tấm ván đặt
trên \VUgras{cái chậu lớn} \textsuperscript{(63)}.

%%K t07-c2-13-1 p
13. --- Cả \VUgras{nhà sữa} \textsuperscript{(59)} được \VUgras{người làm công}
\textsuperscript{(65)} giữ rất sạch; anh vừa rửa xong những \VUgras{bình sữa}
\textsuperscript{(62)} trong chiếc thùng lớn mà anh cầm trên tay.

%%K t07-tit-9 sub
\VUsaut{3.0mm}
\VUpk{10.2pt}{40}{Sân nuôi gia cầm.}
\VUsaut{3.0mm}

%%K t07-c2-14-1 p
14. --- Với đôi guốc gỗ to, anh bước hơi nặng nề, và thế là làm cho con \VUgras{gà tây
trống} \textsuperscript{(68)}, những con \VUgras{vịt mái} \textsuperscript{(66)},
\VUgras{vịt đực} \textsuperscript{(67)} và \VUgras{ngỗng} \textsuperscript{(69)} chạy
tán loạn, chúng há rộng \VUgras{mỏ} \textsuperscript{(70)} và kêu inh ỏi vì sợ.

%%K t07-c2-14-2 p
Con \VUgras{gà trống} \textsuperscript{(71)}, hiếu chiến hơn, dựng chiếc \VUgras{mào}
\textsuperscript{(72)} đỏ, rũ những chiếc lông \VUgras{đuôi} \textsuperscript{(73)} và
cất tiếng « ò ó o » vang dội.

%%K t07-c2-15-1 p
15. --- Một con \VUgras{gà mẹ} \textsuperscript{(74)} bới trên \VUgras{đống phân}
\textsuperscript{(81)} cùng đàn \VUgras{gà con} \textsuperscript{(75)}. Con \VUgras{lợn
con} \textsuperscript{(78)} chạy về phía mẹ nó, con \VUgras{lợn nái}
\textsuperscript{(77)} to lớn mà ta thấy gần chuồng lợn.

%%K t07-c2-16-1 p
16. --- Trong \VUgras{chuồng} của mình, những con \VUgras{thỏ} \textsuperscript{(79)} ăn
ngấu nghiến thứ \VUgras{cỏ} \textsuperscript{(80)} non mà con gái người chủ trại mang
đến. Cô bé Gioanna rất yêu lũ thỏ của mình và mỗi ngày, sau khi đi nhặt \VUgras{trứng}
\textsuperscript{(76)} tươi ở chuồng gà, cô lại đến thăm những người bạn nhỏ ấy.
\VUsaut{6.0mm}
\VUfilet{22mm}
